\section{Preliminaries}

For $I\in\mathcal I(\Theta)$, introduce the three setting-derived chart
pieces
\[
I_0:=I\cap[-1,1],\qquad
I_+:=I\cap(1,\infty),\qquad
I_-:=I\cap(-\infty,-1).
\]
They record the exact interval lengths that appear in the main theorem.  The
corresponding coefficient sweeps are
\[
\begin{aligned}
s_0(\theta;\alpha_{1:d-1})
&:=-\theta^d-\sum_{j=1}^{d-1}\alpha_j\theta^j,
&& |\theta|\leq1,\\
s_\infty(\theta;\alpha_{0:d-2})
&:=-\theta-\sum_{j=0}^{d-2}\alpha_j\theta^{j-d+1},
&& |\theta|\geq1.
\end{aligned}
\]
All empty sums are zero.  These maps express the root equation using
$\alpha_0$ on the inner chart and $\alpha_{d-1}$ on the two outer charts.

The exact chart coefficients and their fixed-$\eta$ envelope are
\[
\begin{aligned}
B_0(d,R)&:=d+\frac{Rd(d-1)}2,\\
B_\infty(d,R)&:=1+\frac{Rd(d-1)}2,\\
M_\eta(d,R)&:=\max\{\bar\kappa_0B_0(d,R),
                         \bar\kappa_\infty B_\infty(d,R)\},\\
\bar\kappa_*&:=\max\{\bar\kappa_0,\bar\kappa_\infty\},\\
P_\eta(d,R)&:=\bar\kappa_*d+\frac{\bar\kappa_*}{2}Rd^2.
\end{aligned}
\]
The main theorem uses $B_0$ and $B_\infty$ for its sharp weighted estimate,
$M_\eta$ for the exact two-chart maximum, and $P_\eta$ for the public
polynomial specialization.  For a scalar map $s$ and a set $J$, the notation
$s(J):=\{s(\theta):\theta\in J\}$ denotes its direct image.
